\section{非线性方程的求解}

\subsubsection{二分法}
\begin{method}[Bisection Method]
	设函数$f$在$[a,b]$上连续，且满足$f(a)f(b)<0$。为求解非线性方程$f(x)=0$，由\cref{lem:IntermediateValueR}可知，存在至少一个根$x^{\star}\in(a,b)$。二分法通过反复对区间进行二等分来逐步缩小根所在的区间：令：
	\begin{equation*}
		c=\frac{a+b}{2}
	\end{equation*}
	若$f(c)=0$，则$c$即为方程的解；否则根据符号判定准则，取满足：
	\begin{equation*}
		f(a)f(c)<0\quad\text{或}\quad f(c)f(b)<0
	\end{equation*}
	的子区间作为新的区间$[a,b]$，并重复上述过程。\par
	经过$n$次迭代后，所得区间长度不超过$(b-a)/2^n$，任意迭代点$x_n$与真根$x^{\star}$的误差满足：
	\begin{equation*}
		|x_n-x^{\star}|\leqslant\frac{b-a}{2^n}
	\end{equation*}
	因此，可通过给定误差容许值$\varepsilon$，选取满足：
	\begin{equation*}
		\frac{b-a}{2^n}\leqslant\varepsilon
	\end{equation*}
	的最小整数$n$来控制计算精度。
\end{method}
\subsubsection{不动点迭代法}
\begin{method}[Fixed-point Iteration Method]
	为求解非线性方程$f(x)=0$，可将其改写为等价的形式$x=T(x)$，若$T$是一个压缩映射，由\cref{theo:ContractionMapTheorem}可求得非线性方程的根。对于迭代次数问题，可用\cref{theo:ContractionMapTheorem}中的第二种误差估计公式来控制精度。
\end{method}
\begin{definition}
	设映射$T$有不动点$x^{\star}$。若存在$x^{\star}$的邻域$\bar{U}(x^{\star},\delta)$满足对任意的$x\in\bar{U}(x^{\star},\delta)$，序列$\{x_n=T^nx\}\to x^{\star}$，则称此时的迭代法\gls{LocalConvergence}。
\end{definition}
\begin{theorem}\label{theo:FixedPointLocalConvergence}
	设映射$T$有不动点$x^{\star}$，$T'$在$x^{\star}$的某个邻域内连续，且$|T'x^{\star}|<1$，则此时的迭代法局部收敛。
\end{theorem}
\begin{proof}
	由条件和\cref{prop:RSeq}(5)可知存在$x^{\star}$的邻域$U$和$L$，使得对任意的$x\in U$有$|T'x|\leqslant L<1$。根据\cref{theo:LagrangeMeanValueTheorem}可知对任意的$x'\in U$有：
	\begin{equation*}
		|Tx-Tx^{\star}|\leqslant L|x-x^{\star}|
	\end{equation*}
	即$T$在$x^{\star}$的某个闭邻域内是压缩映射，由\cref{theo:RnComplete}、\cref{prop:CompleteMetricSpace}(1)和\cref{theo:ContractionMapTheorem}可知此时的迭代法局部收敛。
\end{proof}
\begin{definition}
	设迭代过程$x_{n+1}=Tx_n$收敛于$T$的不动点$x^{\star}$。令$\varepsilon_n=x_n-x^{\star}$，若有：
	\begin{equation*}
		\lim_{n\to+\infty}\frac{\varepsilon_{n+1}}{\varepsilon_n^p}=C\in\mathbb{R}\setminus\{0\}
	\end{equation*}
	则称该迭代过程是$p$阶收敛的。特别的，$p=1$且$|C|<1$时称为\gls{LinearConvergence}，$p>1$时称为\gls{SuperlinearConvergence}。
\end{definition}
\begin{theorem}\label{theo:FixedPointOrderPConvergence}
	设迭代过程$x_{n+1}=Tx_n$。若对于$p\in\mathbb{N}^+$，$T^{(p)}$在$T$的不动点$x^{\star}$的某个邻域内连续，并且有：
	\begin{equation*}
		T'x^{\star}=T''x^{\star}=\cdots=T^{(p-1)}x^{\star}=0,\;T^{(p)}x^{\star}\ne0
	\end{equation*}
	则该迭代过程在$x^{\star}$邻近是$p$阶收敛的。
\end{theorem}
\begin{proof}
	由$T'x^{\star}=0$和\cref{theo:FixedPointLocalConvergence}可知该迭代过程是局部收敛的。根据\info{泰勒展开Lagrange余项}可得：
	\begin{equation*}
		Tx_n=Tx^{\star}+\frac{T^{(p)}\xi_n}{p!}(x_n-x^{\star})^p,\quad\xi_n\text{在$x_n$与$x^{\star}$之间}
	\end{equation*}
	所以：
	\begin{equation*}
		\frac{Tx_n-x^{\star}}{(x_n-x^{\star})^p}=\frac{T^{(p)}\xi_n}{p!}
	\end{equation*}
	由$T^{(p)}$在$x^{\star}$某个邻域内的连续性和\cref{prop:RSeq}(8.c)(4.a)可得：
	\begin{equation*}
		\lim_{n\to+\infty}\frac{Tx_n-x^{\star}}{(x_n-x^{\star})^p}=\lim_{n\to+\infty}\frac{T^{(p)}\xi_n}{p!}=\frac{1}{p!}T^{(p)}\left(\lim_{n\to+\infty}\xi_n\right)=\frac{T^{(p)}x^{\star}}{p!}\qedhere
	\end{equation*}
\end{proof}
\subsubsection{牛顿法}
\begin{method}[Newton Method]
	设$f$在根$x^{\star}$的某个邻域内是$C^2$类函数，且满足$f(x^{\star})=0,\;f'(x^{\star})\ne0$。为求解非线性方程$f(x)=0$，牛顿法对$f$在当前迭代点$x_n$处作一阶Taylor展开，得到：
	\begin{equation*}
		f(x)\approx f(x_n)+f'(x_n)(x-x_n)
	\end{equation*}
	令上式为$0$得到迭代格式为：
	\begin{equation*}
		x_{n+1}=x_n-\frac{f(x_n)}{f'(x_n)}
	\end{equation*}\par
	牛顿法的迭代函数为：
	\begin{equation*}
		Tx=x-\frac{f(x)}{f'(x)}
	\end{equation*}
	于是有：
	\begin{equation*}
		T'x=1-\frac{[f'(x)]^2-f(x)f''(x)}{[f'(x)]^2}=\frac{f(x)f''(x)}{[f'(x)]^2}
	\end{equation*}
	根据\cref{theo:FixedPointOrderPConvergence}可知牛顿法至少是二阶收敛的。
\end{method}